\section{Pułapki i ich unieszkodliwianie}

Jeżeli oddział po raz pierwszy (tylko) wchodzi do jakiejś komnaty jeden z jego członków musi sprawdzić, czy w drzwiach nie została ukryta pułapka.
Podobnie muszą być sprawdzone skrzynie, zawierające skarby.
Pułapka może być unieszkodliwiona przez postać, posiadającą odpowiednie zdolności (,,pułapki'').
Jeżeli próba unieszkodliwienia nie powiedzie się lub nie zostanie podjęta należy rzucić 1K6 i w tabeli 7.1 ,,Pułapki'' sprawdzić rodzaj założonej w danych drzwiach czy skrzyni pułapki.

Gracz, kontrolujący postać, która została wydelegowana do zbadania pułapki rzuca 1K6 i porównuje liczbę wyrzuconych oczek z poziomem zdolności postaci.
Jeżeli wyrzucona liczba oczek jest niższa lub równa poziomowi zdolności pułapka została szczęśliwie unieszkodliwiona, jeżeli wyższa~-- pułapka zadziałała.
Uwaga! Postacie, które nie posiadają odpowiednich zdolności wpadają w pułapki w sposób automatyczny.
Jeżeli pułapka zadziałała gracz, którego postać została wydelegowana do jej zbadania lub gracz, którego postać miała jako pierwsza przejść przez drzwi rzuca 1K6 i w tabeli 7.1 odnajduje typ pułapki, a następnie poniżej (w instrukcji) odczytuje rezultat jej działania.

Tabela ,,Pułapki''~-- patrz plansza z tabelami.

Oddział może się natknąć na 6 typów pułapek:

\begin{enumerate}
    \item
        Strzały~-- ukryty mechanizm wypuszcza strzały.
        Należy rzucić 1K6 i w tabeli 9.9 ,,Rezultat walki'' (kolumna ,,łuk'') odczytać, czy postać została trafiona i ile odniosła ran.
    \item
        Zatrute strzały~-- taka sama procedura, jak przy normalnych strzałach, z tą różnicą, że postać trafiona odnosi dodatkowo 1K3 ran na skutek działania trucizny.
    \item
        Trujący gaz~-- postać, która wpadła w tę pułapkę odnosi 1K3 ran.
    \item
        Wybuch~-- wszystkie postacie, stanowiące oddział (również potwory będące pod wpływem uroku) odnosza 1 ranę.
    \item
        Wrzący olej~-- postać, która wpadła w tę pułapkę odnosi 1K3 ran.
    \item
        Rzuć dwa razy~-- drzwi (skrzynia) zawierają dwie pułapki~-- należy rzucić kostką dwa razy, aby ustalić ich rodzaj.
        Jeżeli znowu zostanie wyrzucona 6-tka należy znów rzucać dwa razy (a więc w sumie trzy pułapki).
        W ten sposób teoretycznie liczba pułapek może być nieskończona.
\end{enumerate}

Poziom zdolności ,,pułapki'' postaci nie może przekroczyć poziomu 5.
